\documentclass[11pt,reqno]{amsart}

\usepackage[T1]{fontenc}
\usepackage{amsmath,amssymb,amsthm}
\usepackage{xcolor}
\usepackage[expansion=false]{microtype}
\usepackage[margin=1.2in]{geometry}
\usepackage[colorlinks=true,linkcolor=blue!50!black,citecolor=blue!50!black,
            urlcolor=blue!50!black]{hyperref}

\theoremstyle{plain}
\newtheorem{theorem}{Theorem}[section]
\newtheorem{proposition}[theorem]{Proposition}
\newtheorem{lemma}[theorem]{Lemma}
\newtheorem{corollary}[theorem]{Corollary}
\theoremstyle{definition}
\newtheorem{remark}[theorem]{Remark}

\newcommand{\R}{\mathbb{R}}
\newcommand{\C}{\mathbb{C}}
\newcommand{\Rp}{\operatorname{Re}}
\newcommand{\Ip}{\operatorname{Im}}
\newcommand{\lxi}[1]{\left(\tfrac{\xi'}{\xi}\right)\!#1}

\title[Shifted screw functions for the Riemann zeta-function]
      {Quantitative asymptotics for Suzuki's shifted screw functions}

\author{Edward B. Baker III}
\email{edwardbaker86@gmail.com}

\date{August 28, 2026}

\begin{document}

\begin{abstract}
Suzuki associated to the Riemann zeta-function a continuous function $\Psi$ and a shifted family
$\Psi_\omega$, and proved that $\zeta(s)\ne0$ for $\Rp s>\tfrac12+\omega$ if and only if
$\Psi_\omega(t)$ is eventually non-negative.  We give a quantitative form of this criterion.
Writing $\Theta=\sup\{\Rp\rho:\xi(\rho)=0\}$ and $\vartheta=\Theta-\tfrac12$, we obtain the exact
decomposition
\[
 \Psi_\omega(t)=\lxi{\left(\tfrac12+\omega\right)}t
 +\lxi{'\left(\tfrac12+\omega\right)}+R_\omega(t),
\]
with an explicit spectral formula for $R_\omega$.  If $\vartheta\le\omega$, then
$|R_\omega(t)|\le A^*(\omega)e^{-(\omega-\vartheta)t}$ with an explicit constant $A^*(\omega)$;
thus every valid shifted criterion has an
explicit positive linear margin.  If instead a zero lies to the right of the shifted boundary, its
contribution grows exponentially at a rate equal to its horizontal excess.  A Landau-type argument
then gives a sharp dichotomy across the critical shift and several one-sided consequences, including
an equivalent formulation of RH by integrability of the negative part of $\Psi$.  These results
quantify Suzuki's criterion; they do not by themselves provide new zero-free information.
\end{abstract}

\maketitle

% ==========================================================================
\section{Introduction}
% ==========================================================================

Let
\[
 \xi(s)=\tfrac12 s(s-1)\pi^{-s/2}\Gamma\!\left(\tfrac s2\right)\zeta(s),
\]
so that $\xi$ is entire of order one and satisfies $\xi(s)=\xi(1-s)$.  Put
$\Xi(z)=\xi(\tfrac12-iz)$, and denote the zeros of $\Xi$ by $\gamma$.  Thus, if
$\rho=\beta+i\tau$ is a zero of $\xi$, then
\begin{equation}\label{eq:gammarho}
 \gamma=i\left(\rho-\tfrac12\right)=-\tau+i\left(\beta-\tfrac12\right).
\end{equation}
The Riemann Hypothesis (RH) is therefore equivalent to the reality of every $\gamma$.

Suzuki \cite{Suzuki2023} introduced the continuous even function
\begin{equation}\label{eq:PsiTwoFaces}
 \Psi(t)=\sum_\gamma\frac{1-\cos(\gamma t)}{\gamma^2},
 \qquad
 \int_0^\infty \Psi(t)e^{izt}\,dt
 =-\frac1{z^2}\lxi{\left(\tfrac12-iz\right)}
 \quad (\Ip z>\tfrac12),
\end{equation}
and proved
\begin{equation}\label{eq:SuzukiRH}
 \mathrm{RH}\quad\Longleftrightarrow\quad \Psi(t)\ge0\quad\text{for every }t.
\end{equation}
He also introduced, for $\omega\in\R$, the shifted family
\begin{equation}\label{eq:Psiomegadef}
 \Psi_\omega(t)
 :=e^{-\omega t}\Psi(t)
 +2\omega\int_0^t e^{-\omega u}\Psi(u)\,du
 +\omega^2\int_0^t(t-u)e^{-\omega u}\Psi(u)\,du,
 \qquad t>0,
\end{equation}
extended evenly, for which
\begin{equation}\label{eq:PsiomegaLaplace}
 \int_0^\infty\Psi_\omega(t)e^{izt}\,dt
 =-\frac1{z^2}\lxi{\left(\tfrac12+\omega-iz\right)},
 \qquad \Ip z>\tfrac12-\omega.
\end{equation}
His shifted criterion is
\begin{equation}\label{eq:SuzukiShifted}
 \xi(s)\ne0\ \text{for }\Rp s>\tfrac12+\omega
 \quad\Longleftrightarrow\quad
 \Psi_\omega(t)\ge0\ \text{for all sufficiently large }t.
\end{equation}
Moreover, in the forward direction Suzuki's Herglotz argument, through the Kre\u\i n--Langer
theory of screw functions \cite{KreinLanger}, gives pointwise non-negativity for
all $t$; see \cite[\S11]{Suzuki2023}.  The shifted family also satisfies a composition law, so
non-negativity propagates upward in $\omega$.

The purpose of this note is to extract the quantitative information already latent in Suzuki's
spectral formula.  Set
\[
 \Theta:=\sup\{\Rp\rho:\xi(\rho)=0\},
 \qquad \vartheta:=\Theta-\tfrac12\in[0,\tfrac12].
\]
By the functional equation, every zero $\gamma=a+ib$ of $\Xi$ satisfies $|b|\le\vartheta$.
The main observation is that $\Psi_\omega$ splits into an unconditional linear term and an
oscillatory remainder.  The slope and intercept of the linear term are the classical quantities
$(\xi'/\xi)(\tfrac12+\omega)$ and $(\xi'/\xi)'(\tfrac12+\omega)$, while the remainder records the
horizontal location of the zeros through factors $e^{(\Rp\rho-1/2-\omega)t}$.  This yields three
consequences.  First, whenever the half-plane in \eqref{eq:SuzukiShifted} is zero-free, the
criterion is not asymptotically saturated: it has a positive linear margin with an explicit onset.
Second, a zero beyond the shifted boundary produces an exponentially growing oscillation whose
rate is exactly its horizontal excess.  Third, the same Laplace-transform argument underlying
Suzuki's criterion converts one-sided growth estimates, including averaged estimates for the
negative part, into zero-free half-planes.

The statements below are quantitative reformulations of a known zero-free criterion.  In
particular, proving the needed one-sided estimates from the arithmetic side would itself amount to
proving new zero-free information.  No such estimate is claimed here.

The functions $\Psi$ and $\Psi_\omega$, the transform \eqref{eq:PsiomegaLaplace}, the criterion
\eqref{eq:SuzukiShifted}, and the Herglotz/screw-function interpretation are due to Suzuki
\cite{Suzuki2023}.  The use of Herglotz positivity at a general abscissa is closely related to
Lagarias \cite{Lagarias1999}.  The contributions here begin with Proposition \ref{prop:AB} and
Theorem \ref{thm:asym}.

% ==========================================================================
\section{Exact decomposition of the shifted family}\label{sec:decomposition}
% ==========================================================================

We use the standard Hadamard consequences
\begin{equation}\label{eq:hadamard}
 \frac{\Xi'}{\Xi}(z)=\sum_\gamma\frac1{z-\gamma},
 \qquad
 \frac{\Xi'}{\Xi}(z)=-i\,\lxi{\left(\tfrac12-iz\right)},
 \qquad
 \lxi{(1-s)}=-\lxi{(s)},
\end{equation}
where the first sum is taken with the usual pairing of $\gamma$ and $-\gamma$; see, for example,
\cite[\S12]{Davenport}.  Suzuki's series for the shifted family is
\begin{equation}\label{eq:Psiomegaseries}
 \Psi_\omega(t)=\sum_\gamma\frac{1}{(\gamma^2+\omega^2)^2}
 \Bigl[\omega t(\gamma^2+\omega^2)+(\gamma^2-\omega^2)
 -e^{-\omega t}(\gamma^2-\omega^2)\cos(\gamma t)
 -2\gamma\omega e^{-\omega t}\sin(\gamma t)\Bigr].
\end{equation}

For $0<\omega<\tfrac12$ define
\begin{equation}\label{eq:AB}
 A(\omega):=\sum_\gamma\frac1{\gamma^2+\omega^2},
 \qquad
 B(\omega):=\sum_\gamma\frac{\gamma^2-\omega^2}{(\gamma^2+\omega^2)^2},
\end{equation}
and
\begin{equation}\label{eq:Astar}
 A^*(\omega):=\frac12\sum_\gamma
 \left(\frac1{|\gamma-i\omega|^2}+\frac1{|\gamma+i\omega|^2}\right).
\end{equation}
At $\omega=0$ we use the evident limiting definition for $A^*(0)$.

\begin{lemma}\label{lem:conv}
For $0<\omega<\tfrac12$, the series in \eqref{eq:AB} and \eqref{eq:Astar} converge absolutely, and
\[
 \sum_\gamma |\gamma|\,|\gamma^2+\omega^2|^{-2}<\infty.
\]
\end{lemma}

\begin{proof}
Write $\gamma=a+ib$.  Unconditionally $|b|\le\tfrac12$, and the least positive ordinate of a
nontrivial zero of $\zeta$ is $14.1347\ldots$, so $|a|\ge14$; see \cite{Titchmarsh}.  Hence
\[
 \Rp(\gamma^2+\omega^2)=a^2-b^2+\omega^2\ge a^2-\tfrac14>0,
 \qquad
 |\gamma\pm i\omega|^2\ge a^2.
\]
The Riemann--von Mangoldt formula then gives the asserted convergence by partial summation.
\end{proof}

\begin{proposition}\label{prop:AB}
For $0<\omega<\tfrac12$,
\begin{equation}\label{eq:ABclosed}
 A(\omega)=\frac1\omega\lxi{\left(\tfrac12+\omega\right)},
 \qquad
 B(\omega)=\lxi{'\left(\tfrac12+\omega\right)}.
\end{equation}
Moreover $A$ is strictly decreasing on $(0,\tfrac12)$ and
\begin{equation}\label{eq:Aendpoint}
 \lim_{\omega\to1/2^-}A(\omega)
 =\sum_\rho\frac1{\rho(1-\rho)}
 =2+\gamma_0-\log 4\pi
 =0.0461914\ldots.
\end{equation}
If RH holds, then $A^*(\omega)=A(\omega)$.
\end{proposition}

\begin{proof}
Evaluating \eqref{eq:hadamard} at $z=\pm i\omega$ gives
\[
 \sum_\gamma\frac1{\gamma-i\omega}
 =i\lxi{\left(\tfrac12+\omega\right)},
 \qquad
 \sum_\gamma\frac1{\gamma+i\omega}
 =-i\lxi{\left(\tfrac12+\omega\right)}.
\]
Partial fractions and absolute convergence therefore give the first identity in
\eqref{eq:ABclosed}.  Since
\[
 \frac{\partial}{\partial\omega}\frac{\omega}{\gamma^2+\omega^2}
 =\frac{\gamma^2-\omega^2}{(\gamma^2+\omega^2)^2},
\]
local uniform convergence permits termwise differentiation, yielding the second identity.

For monotonicity, note first that a real $\gamma=a$ contributes $(a^2+\omega^2)^{-1}$, which is
strictly decreasing in $\omega$.  For the non-real zeros, pair $\gamma=a+ib$ with
$\overline\gamma$.  Their contribution to $A$ is
\[
 \frac{2X}{X^2+Y^2},
 \qquad
 X=a^2-b^2+\omega^2,
 \quad Y=2ab.
\]
Differentiating with respect to $\omega^2$ gives
\[
 \frac{2(Y^2-X^2)}{(X^2+Y^2)^2}<0,
\]
because
\[
 X-|Y|=a^2-b^2+\omega^2-2|ab|
 \ge 14^2-\tfrac14-14>0.
\]
Thus every conjugate pair decreases strictly.  At $\omega=\tfrac12$,
$\gamma^2+\tfrac14=\rho(1-\rho)$ by \eqref{eq:gammarho}, and the classical identity in
\eqref{eq:Aendpoint} follows from the Hadamard product; dominated convergence gives the limit.
Under RH all $\gamma$ are real, and each summand in \eqref{eq:Astar} equals
$(\gamma^2+\omega^2)^{-1}$ after averaging the two signs.
\end{proof}

\begin{remark}\label{rem:B0}
The functional equation makes $\xi'/\xi$ odd about $s=\tfrac12$.  Hence
\begin{equation}\label{eq:B0}
 B(0):=\left(\frac{\xi'}{\xi}\right)'\!\left(\tfrac12\right)
 =\sum_\gamma\frac1{\gamma^2}
 =0.0462099862308\ldots,
\end{equation}
and
\begin{equation}\label{eq:slopeSmallOmega}
 \lxi{\left(\tfrac12+\omega\right)}=B(0)\,\omega+O(\omega^3).
\end{equation}
\end{remark}

\begin{theorem}[Exact decomposition and zero-free asymptotic]\label{thm:asym}
Let $0<\omega<\tfrac12$.  For every $t\ge0$,
\begin{equation}\label{eq:mainDecomp}
 \Psi_\omega(t)=\lxi{\left(\tfrac12+\omega\right)}t
 +\lxi{'\left(\tfrac12+\omega\right)}+R_\omega(t),
\end{equation}
where
\begin{equation}\label{eq:Rclean}
 R_\omega(t)=-\frac{e^{-\omega t}}2\sum_\gamma
 \left[
 \frac{e^{i\gamma t}}{(\gamma+i\omega)^2}
 +\frac{e^{-i\gamma t}}{(\gamma-i\omega)^2}
 \right].
\end{equation}
Consequently, unconditionally,
\begin{equation}\label{eq:Rtrivial}
 |R_\omega(t)|\le A^*(\omega)e^{(1/2-\omega)t}.
\end{equation}
More generally, if $\vartheta\le\omega$, then
\begin{equation}\label{eq:RzeroFree}
 |R_\omega(t)|\le A^*(\omega)e^{-(\omega-\vartheta)t}.
\end{equation}
Under RH this becomes
\begin{equation}\label{eq:RRH}
 |R_\omega(t)|\le A(\omega)e^{-\omega t}
 =\frac1\omega\lxi{\left(\tfrac12+\omega\right)}e^{-\omega t}.
\end{equation}
\end{theorem}

\begin{proof}
The first two terms of \eqref{eq:Psiomegaseries} sum to
\[
 \omega A(\omega)t+B(\omega)
 =\lxi{\left(\tfrac12+\omega\right)}t
 +\lxi{'\left(\tfrac12+\omega\right)}
\]
by Proposition \ref{prop:AB}.  For the remaining terms, the identity
\[
 (\gamma^2-\omega^2)\cos(\gamma t)+2\gamma\omega\sin(\gamma t)
 =\tfrac12e^{i\gamma t}(\gamma-i\omega)^2
 +\tfrac12e^{-i\gamma t}(\gamma+i\omega)^2
\]
gives \eqref{eq:Rclean} after division by
$(\gamma^2+\omega^2)^2=(\gamma-i\omega)^2(\gamma+i\omega)^2$.
If $\gamma=a+ib$, then $|e^{\pm i\gamma t}|=e^{\mp bt}$.  Using $|b|\le\tfrac12$ gives
\eqref{eq:Rtrivial}, while $|b|\le\vartheta$ gives \eqref{eq:RzeroFree}.  Under RH,
$A^*=A$ by Proposition \ref{prop:AB}.
\end{proof}

% ==========================================================================
\section{The margin and the critical shift}\label{sec:margin}
% ==========================================================================

\begin{corollary}[Linear margin]\label{cor:margin}
For $0<\omega<\tfrac12$, the slope in \eqref{eq:mainDecomp} is positive and satisfies
\begin{equation}\label{eq:slopeLower}
 \lxi{\left(\tfrac12+\omega\right)}
 =\omega A(\omega)
 \ge (2+\gamma_0-\log4\pi)\,\omega.
\end{equation}
It tends to zero linearly as $\omega\to0^+$ according to \eqref{eq:slopeSmallOmega}.
\end{corollary}

\begin{proof}
This is Proposition \ref{prop:AB}, Remark \ref{rem:B0}, and the monotonicity of $A$.
\end{proof}

Thus the shifted criterion is qualitatively different from the unshifted one at infinity: for every
$\omega>0$ satisfying the corresponding zero-free hypothesis, the leading term grows linearly.
The next statement makes this quantitative.

\begin{theorem}[Explicit onset of the margin]\label{thm:threshold}
Let $0<\omega<\tfrac12$ and assume $\vartheta\le\omega$.  Then
\begin{equation}\label{eq:lowerMargin}
 \Psi_\omega(t)\ge
 \lxi{\left(\tfrac12+\omega\right)}t
 +\lxi{'\left(\tfrac12+\omega\right)}-A^*(\omega)
 \qquad(t\ge0).
\end{equation}
In particular $\Psi_\omega(t)>0$ for
\begin{equation}\label{eq:t0}
 t>t_0(\omega):=
 \frac{A^*(\omega)-B(\omega)}{\omega A(\omega)}.
\end{equation}
Moreover
\begin{equation}\label{eq:AstarMinusBbound}
 A^*(\omega)-B(\omega)
 \le8\omega^2\sum_\rho\frac1{(\Ip\rho)^4},
\end{equation}
and hence
\begin{equation}\label{eq:t0numeric}
 t_0(\omega)\le0.0129\,\omega<0.0065.
\end{equation}
\end{theorem}

\begin{proof}
The bound \eqref{eq:lowerMargin} follows immediately from Theorem \ref{thm:asym} because
$e^{-(\omega-\vartheta)t}\le1$.  To estimate the difference between the intercept and the
oscillatory amplitude, write $\gamma=a+ib$ and $z_1=\gamma-i\omega$, $z_2=\gamma+i\omega$.
Since
\[
 \frac{\gamma^2-\omega^2}{(\gamma^2+\omega^2)^2}
 =\frac12\left(\frac1{z_1^2}+\frac1{z_2^2}\right)
\]
and, for $z=x+iy\ne0$,
\[
 \frac1{|z|^2}-\Rp\frac1{z^2}=\frac{2y^2}{|z|^4},
\]
summing over conjugate pairs gives the exact identity
\begin{equation}\label{eq:AstarminusB}
 A^*(\omega)-B(\omega)
 =\sum_\gamma\left[
 \frac{(b-\omega)^2}{|\gamma-i\omega|^4}
 +\frac{(b+\omega)^2}{|\gamma+i\omega|^4}
 \right].
\end{equation}
Under $|b|\le\vartheta\le\omega$, each numerator is at most $4\omega^2$ and
$|\gamma\pm i\omega|^2\ge a^2$, proving \eqref{eq:AstarMinusBbound}.  The numerical value
\[
 \sum_\rho(\Ip\rho)^{-4}=7.4345\ldots\times10^{-5}<7.44\times10^{-5}
\]
may be obtained from the first ordinates together with a Riemann--von Mangoldt tail estimate; an
explicit error bound such as \cite{TrudgianArg} is more than sufficient at the displayed precision.
Combining this with \eqref{eq:slopeLower} gives
$t_0(\omega)\le 8\cdot7.44\cdot10^{-5}\,(2+\gamma_0-\log4\pi)^{-1}\omega=0.01288\ldots\times\omega$,
whence \eqref{eq:t0numeric}.
\end{proof}

\begin{remark}
Under $\vartheta\le\omega$, Suzuki's Herglotz argument already gives $\Psi_\omega\ge0$ for every
$t$ (\S1).  What Theorem \ref{thm:threshold} adds is the linear lower bound
\eqref{eq:lowerMargin} --- wherever the shifted criterion holds, it holds with a margin that grows
without bound, so it is never asymptotically saturated --- together with an elementary proof that
avoids the Kre\u\i n--Langer correspondence.
\end{remark}

The qualitative boundary is the critical shift $\omega_c:=\vartheta=\Theta-\tfrac12$.
At and above it the remainder in \eqref{eq:mainDecomp} is bounded; below it, one-sided bounds of any
smaller exponential type are impossible.

\begin{theorem}[Critical-shift dichotomy]\label{thm:transition}
Let $0<\omega<\tfrac12$.
\begin{enumerate}
\item[(i)] If $\vartheta\le\omega$, then
\[
 \Psi_\omega(t)=\lxi{\left(\tfrac12+\omega\right)}t+O(1),
 \qquad
 \frac{\Psi_\omega(t)}t\longrightarrow
 \lxi{\left(\tfrac12+\omega\right)}.
\]
\item[(ii)] If $\omega<\vartheta$, then for every $0\le\varepsilon<\vartheta-\omega$,
\[
 \limsup_{t\to\infty}e^{-\varepsilon t}\Psi_\omega(t)=+\infty,
 \qquad
 \liminf_{t\to\infty}e^{-\varepsilon t}\Psi_\omega(t)=-\infty.
\]
\end{enumerate}
\end{theorem}

\begin{proof}
Part (i) follows from Theorem \ref{thm:asym}.  For (ii), suppose for contradiction that
\[
 \Psi_\omega(t)\ge-Ce^{\varepsilon t}-D
\]
for some $C,D>0$ and $0\le\varepsilon<\vartheta-\omega$.  Set
$f(t)=\Psi_\omega(t)+Ce^{\varepsilon t}+D\ge0$.  From \eqref{eq:Rclean},
\[
 |\Psi_\omega(t)|\le \omega A(\omega)t+|B(\omega)|+A^*(\omega)e^{(\vartheta-\omega)t},
\]
so the Laplace transform $F(z)=\int_0^\infty f(t)e^{izt}\,dt$ converges for
$\Ip z>\vartheta-\omega$.  On that half-plane, by analytic continuation from
$\Ip z>\tfrac12-\omega$,
\begin{equation}\label{eq:Ftransform}
 F(z)=-\frac1{z^2}\lxi{\left(\tfrac12+\omega-iz\right)}
 +\frac{iC}{z-i\varepsilon}+\frac{iD}{z}.
\end{equation}
Let $\sigma_c$ be its abscissa of convergence.  Since $f\ge0$, Landau's theorem for Laplace
transforms \cite[Ch.~II, Thm.~5b]{Widder} implies that $i\sigma_c$ is a singularity of $F$.
On the ray $\{iy:y>\varepsilon\}$, however, the right-hand side of \eqref{eq:Ftransform} can have
a singularity only if $\tfrac12+\omega+y$ is a real zero of $\xi$.  There are no such zeros in
$(\tfrac12,1]$.  Hence $\sigma_c\le\varepsilon$.  Analytic continuation of
\eqref{eq:Ftransform} to $\Ip z>\varepsilon$ then excludes every zero with
$\Rp\rho>\tfrac12+\omega+\varepsilon$, contradicting
$\omega+\varepsilon<\vartheta$.  This proves the $\liminf$ statement.  Applying the same argument
to $Ce^{\varepsilon t}+D-\Psi_\omega(t)$ proves the $\limsup$ statement.
\end{proof}

\begin{corollary}\label{cor:boundedCharacterization}
For $0<\omega<\tfrac12$,
\begin{equation}\label{eq:boundedCharacterization}
 \Theta\le\tfrac12+\omega
 \quad\Longleftrightarrow\quad
 \Psi_\omega(t)-\lxi{\left(\tfrac12+\omega\right)}t=O(1).
\end{equation}
Consequently
\[
 \Theta-\tfrac12
 =\inf\left\{\omega>0:
 \Psi_\omega(t)-\lxi{\left(\tfrac12+\omega\right)}t\ \text{is bounded}\right\}.
\]
\end{corollary}

\begin{proof}
The forward direction is Theorem \ref{thm:transition}(i).  Conversely, boundedness contradicts
Theorem \ref{thm:transition}(ii) for every $\varepsilon>0$ if $\omega<\vartheta$.
\end{proof}

The spectral meaning of the exponent in Theorem \ref{thm:transition}(ii) can be seen directly from
a single zero quartet.

\begin{proposition}[Contribution of a zero beyond the shifted boundary]\label{prop:sens}
Let $\rho_0=\tfrac12+\delta+i\tau_0$ be a zero of $\xi$ with $\delta>\omega\ge0$ and $\tau_0>0$.
Let $\Psi_\omega^{(\rho_0)}$ denote the contribution to Suzuki's spectral series of the quartet
$\{\rho_0,\overline{\rho_0},1-\rho_0,1-\overline{\rho_0}\}$.  Then
\begin{equation}\label{eq:singleZero}
 \Psi_\omega^{(\rho_0)}(t)
 =-\frac{2e^{(\delta-\omega)t}}{\tau_0^2+(\delta-\omega)^2}
 \cos(\tau_0t+\alpha)+O(t),
 \qquad
 \alpha=2\arg\bigl(\tau_0+i(\delta-\omega)\bigr).
\end{equation}
In particular,
\[
 \limsup_{t\to\infty}e^{-(\delta-\omega)t}
 \left|\Psi_\omega^{(\rho_0)}(t)\right|
 =\frac{2}{\tau_0^2+(\delta-\omega)^2}.
\]
\end{proposition}

\begin{proof}
By \eqref{eq:gammarho}, the quartet corresponds to
$\gamma\in\{\pm\tau_0\pm i\delta\}$.  For $\gamma=a+i\delta$ with $a=\pm\tau_0$, the dominant term
in \eqref{eq:Rclean} is
\[
 -\frac{e^{(\delta-\omega)t}e^{-iat}}
 {2\bigl(a+i(\delta-\omega)\bigr)^2}
 \bigl(1+O(e^{-2\delta t})\bigr).
\]
Summing the four dominant contributions gives the oscillatory term in \eqref{eq:singleZero}; the
non-oscillatory contribution of the quartet to the first two terms of \eqref{eq:mainDecomp} is
$O(t)$.
\end{proof}

% ==========================================================================
\section{One-sided consequences}\label{sec:onesided}
% ==========================================================================

The proof of Theorem \ref{thm:transition} uses only one-sided control.  We record three consequences
because they give compact equivalents of the same zero-free information.  They should be understood
as Laplace-transform reformulations, not as independent evidence for RH; in substance they are
twice-integrated, damped analogues of the classical one-sided oscillation theorems of E.~Schmidt
and Landau (see \cite[Ch.~V]{Ingham}), with the screw normalisation supplying exact constants and
a single continuous carrier function.

\begin{proposition}[One further integration]\label{prop:integrated}
Define
\[
 \Psi^{[1]}(t):=\int_0^t\Psi(u)\,du.
\]
Then RH holds if and only if $\Psi^{[1]}(t)\ge0$ for all sufficiently large $t$.  Under RH,
\begin{equation}\label{eq:integratedRH}
 \Psi^{[1]}(t)=B(0)t-\sum_\gamma\frac{\sin(\gamma t)}{\gamma^3},
 \qquad
 \left|\sum_\gamma\frac{\sin(\gamma t)}{\gamma^3}\right|
 \le \sum_\gamma|\gamma|^{-3}=0.0014591\ldots,
\end{equation}
so $\Psi^{[1]}(t)>0$ for $t>0.0316$.
\end{proposition}

\begin{proof}
Under RH, the series in \eqref{eq:PsiTwoFaces} converges uniformly, and termwise integration gives
\eqref{eq:integratedRH}.  Thus eventual positivity is necessary.  Unconditionally
$\Psi^{[1]}(t)=O(e^{t/2})$, and integration by parts in \eqref{eq:PsiTwoFaces} gives
\begin{equation}\label{eq:integratedTransform}
 \int_0^\infty\Psi^{[1]}(t)e^{izt}\,dt
 =\frac1{iz^3}\lxi{\left(\tfrac12-iz\right)},
 \qquad \Ip z>\tfrac12.
\end{equation}
If $\Psi^{[1]}$ is eventually non-negative, add a constant so that it is non-negative on all of
$[0,\infty)$.  Landau's theorem then places a singularity of its transform on the positive
imaginary axis at the abscissa of convergence.  The meromorphic continuation in
\eqref{eq:integratedTransform} has no such singularity corresponding to a real zero in
$(\tfrac12,1]$.  The same argument as in Theorem \ref{thm:transition} therefore forces the
abscissa to be at most zero and excludes zeros with real part $>\tfrac12$.
\end{proof}

\begin{corollary}[Pointwise one-sided bounds]\label{cor:oneside}
Let $0\le\omega<\tfrac12$ and $\varepsilon\ge0$.
If, for some $C>0$, either
\[
 \Psi_\omega(t)\ge-Ce^{\varepsilon t}
 \quad\text{or}\quad
 \Psi_\omega(t)\le Ce^{\varepsilon t}
\]
for all sufficiently large $t$, then
\begin{equation}\label{eq:onesideConclusion}
 \Theta\le\tfrac12+\omega+\varepsilon.
\end{equation}
In particular,
\begin{equation}\label{eq:RHboundedBelow}
 \mathrm{RH}\quad\Longleftrightarrow\quad
 \liminf_{t\to\infty}\Psi(t)>-\infty.
\end{equation}
Also, if for one $\omega\in[0,\tfrac12)$ and some $\delta\in(0,\tfrac12-\omega]$ one proves
\[
 \Psi_\omega(t)\ge-Ce^{(1/2-\omega-\delta)t}
\]
eventually, then $\Theta\le1-\delta$.
\end{corollary}

\begin{proof}
The proof of Theorem \ref{thm:transition}(ii), and its sign-reversed version, applies verbatim with
$\varepsilon$ in place of a number smaller than $\vartheta-\omega$.  If
$\Theta>\tfrac12+\omega+\varepsilon$, choose
$\varepsilon'<\vartheta-\omega$ with $\varepsilon<\varepsilon'$ and obtain a contradiction to the
assumed one-sided bound.  The case $\omega=0$ follows from the same transform argument using
\eqref{eq:PsiTwoFaces}.
\end{proof}

Write $\Psi_\omega^-(t)=\max\{-\Psi_\omega(t),0\}$.

\begin{theorem}[Averaged one-sided form]\label{thm:L1}
Let $0\le\omega<\tfrac12$ and $\varepsilon\ge0$.  If
\begin{equation}\label{eq:L1hyp}
 \int_0^T\Psi_\omega^-(t)\,dt\le Ce^{\varepsilon T}
 \qquad(T\ge0),
\end{equation}
then
\begin{equation}\label{eq:L1conclusion}
 \Theta\le\tfrac12+\omega+\varepsilon.
\end{equation}
Conversely, if $\vartheta\le\omega+\varepsilon$, then \eqref{eq:L1hyp} holds for some $C$.
In particular,
\begin{equation}\label{eq:RHnegativePart}
 \mathrm{RH}
 \quad\Longleftrightarrow\quad
 \int_0^\infty\Psi^-(t)\,dt<\infty.
\end{equation}
\end{theorem}

\begin{proof}
Set $g=\Psi_\omega^+=\Psi_\omega+\Psi_\omega^-\ge0$ and
\[
 h(z)=\int_0^\infty\Psi_\omega^-(t)e^{izt}\,dt.
\]
If $H(T)=\int_0^T\Psi_\omega^-(t)\,dt$, then integration by parts gives, for
$\Ip z>\varepsilon$, absolute convergence of $h$ from $H(T)\le Ce^{\varepsilon T}$; local uniform
convergence gives analyticity.  On $\Ip z>\tfrac12-\omega$,
\begin{equation}\label{eq:gTransform}
 \widehat g(z)
 =-\frac1{z^2}\lxi{\left(\tfrac12+\omega-iz\right)}+h(z).
\end{equation}
Since $g\ge0$, Landau's theorem makes $i\sigma_c$ a singularity of $\widehat g$, where $\sigma_c$
is its abscissa of convergence.  If $\sigma_c>\varepsilon$, then $h$ is analytic near
$i\sigma_c$, while the first term on the right of \eqref{eq:gTransform} can be singular there only
if $\tfrac12+\omega+\sigma_c$ is a real zero of $\xi$ in $(\tfrac12,1]$.  There are no such zeros.
Hence $\sigma_c\le\varepsilon$.  It follows from \eqref{eq:gTransform} that
$-z^{-2}(\xi'/\xi)(\tfrac12+\omega-iz)$ is analytic on $\Ip z>\varepsilon$, proving
\eqref{eq:L1conclusion}.

Conversely, if $\vartheta\le\omega$, Theorem \ref{thm:threshold} shows that $\Psi_\omega^-$ has
bounded support (and for $\omega=0$, RH gives $\Psi^-=0$).  If
$\omega<\vartheta\le\omega+\varepsilon$, Theorem \ref{thm:asym} gives
\[
 \Psi_\omega^-(t)
 \le A^*(\omega)e^{(\vartheta-\omega)t}+|B(\omega)|,
\]
whose integral over $[0,T]$ is $O(e^{\varepsilon T})$.
\end{proof}

% ==========================================================================
\section{Final remarks}
% ==========================================================================

The decomposition \eqref{eq:mainDecomp} makes Suzuki's shifted criterion a competition between an
explicit linear term and an oscillatory spectral remainder.  At or to the right of the critical
shift, the remainder is bounded and the linear term wins.  To the left, a zero with
$\Rp\rho>\tfrac12+\omega$ contributes an exponential factor and ultimately defeats every linear
margin.  The shifted criterion therefore transports the zero-free-half-plane problem faithfully;
it does not remove its analytic difficulty.

The same observation explains why the one-sided formulations in Section \ref{sec:onesided} should
not be interpreted as progress toward RH by themselves.  For example, obtaining any bound of the
form
\[
 \Psi_\omega(t)\ge-Ce^{(1/2-\omega-\delta)t}
\]
with $0<\delta\le\tfrac12-\omega$ would already prove the zero-free half-plane $\Theta\le1-\delta$.  The value of the
reformulation is instead that the required information can be expressed as a one-sided estimate for
a single continuous function, or even as an averaged estimate for its negative part.

Two extensions seem natural but are not needed for the core statements above.  First, verified
zeros together with explicit zero-free regions can be inserted into \eqref{eq:Rclean} to produce
large unconditional initial intervals on which $\Psi_\omega$ is positive.  Making the resulting
numerical constants fully rigorous requires a careful Stieltjes treatment of the zero-counting
error and is best separated from the present argument.  Second, analogous formulas hold for
self-dual Dirichlet $L$-functions; off-central real zeros behave differently because their
contribution has no oscillatory phase.  Both directions may be worth developing after the basic
shifted asymptotic has been independently checked.

\begin{thebibliography}{99}

\bibitem{Davenport} H.~Davenport,
\emph{Multiplicative Number Theory}, 3rd ed., Graduate Texts in Mathematics \textbf{74},
Springer, New York, 2000.

\bibitem{Ingham} A.~E.~Ingham,
\emph{The Distribution of Prime Numbers}, Cambridge Tracts in Mathematics \textbf{30},
Cambridge Univ. Press, Cambridge, 1932.

\bibitem{KreinLanger} M.~G.~Kre\u\i n and H.~Langer,
\emph{Continuation of hermitian positive definite functions and related questions},
Integral Equations Operator Theory \textbf{78} (2014), 1--69.

\bibitem{Lagarias1999} J.~C.~Lagarias,
\emph{On a positivity property of the Riemann $\xi$-function},
Acta Arith. \textbf{89} (1999), 217--234.

\bibitem{Suzuki2023} M.~Suzuki,
\emph{Aspects of the screw function corresponding to the Riemann zeta-function},
J. London Math. Soc. (2) \textbf{108} (2023), 1448--1487; arXiv:2206.03682.

\bibitem{Titchmarsh} E.~C.~Titchmarsh,
\emph{The Theory of the Riemann Zeta-Function}, 2nd ed., revised by D.~R.~Heath-Brown,
Oxford Univ. Press, 1986.

\bibitem{TrudgianArg} T.~S.~Trudgian,
\emph{An improved upper bound for the argument of the Riemann zeta-function on the critical line II},
J. Number Theory \textbf{134} (2014), 280--292.

\bibitem{Widder} D.~V.~Widder,
\emph{The Laplace Transform}, Princeton Math. Series \textbf{6}, Princeton Univ. Press, 1941.

\end{thebibliography}

\end{document}
